% Auto-generated by generate_tex_fragments.py -- do not hand-edit.
\begin{longtable}{@{}L{3.0cm}L{2.4cm}L{9.0cm}@{}}
\caption{Evidence gaps, excerpted from the gaps\_noted field logged in each search-*.json angle file at the time of the search pass. Text is the search log's own statement, condensed to fit the table column, not a new claim.}\label{tab:gaps}\\
\toprule
\textbf{Gap theme} & \textbf{Search angle} & \textbf{Logged gap statement}\\
\midrule
\endfirsthead
\multicolumn{3}{l}{\textit{Table \thetable{} continued}}\\
\toprule
\textbf{Gap theme} & \textbf{Search angle} & \textbf{Logged gap statement}\\
\midrule
\endhead
\bottomrule
\endfoot
Thai \slash{} Thai-Muslim samples & assessor-\seqsplit{independence}-checkin-models & No Thailand-specific or premarital-education-specific record was found for this angle; angle (a)-(d) evidence is drawn from adjacent fields (organ \seqsplit{transplantation} ethics, IPV risk assessment, suicide\slash{}self-harm brief-contact trials, US pediatric\slash{}prevention-science \seqsplit{implementation} studies) and must be explicitly flagged as analogical transfer, not direct T-PHE evidence, in any manuscript use.\\
\\[3pt]
Thai \slash{} Thai-Muslim samples & couple-work-safety-coercive-control & No PubMed-indexed study was found evaluating a couple-based premarital or pre-nikah education programme with an embedded coercive-control screening\slash{}stop-rule mechanism in a Muslim or Thai\slash{}Southeast Asian \seqsplit{intercultural} setting --- the exact population T-PHE targets.\\
\\[3pt]
Thai \slash{} Thai-Muslim samples & decision-autonomy-measures & No PubMed-indexed decision-autonomy or coercion-\seqsplit{recognition} instrument was found validated \seqsplit{specifically} in a Thai population, Thai Muslim population, or Thai \seqsplit{intercultural}-couple population; the closest cultural\slash{}religious matches found (Türkiye, Lebanon, Rohingya\slash{}Syrian refugees) are \seqsplit{geographically} and culturally distant from Thailand.\\
\\[3pt]
Thai \slash{} Thai-Muslim samples & decision-autonomy-measures & Cross-cultural adaptation studies located here show mixed results (Hazar 2024 Turkey: acceptable but uneven subscale \seqsplit{reliability}; Dagher 2024 Lebanon: sub-optimal CFA fit even after item removal) --- this is direct evidence against assuming any existing decision-control\slash{}autonomy instrument would transfer cleanly to a new Thai Muslim\slash{}\seqsplit{intercultural} context without \seqsplit{independent} local validation, a gap the T-PHE abstract...\\
\\[3pt]
Thai \slash{} Thai-Muslim samples & digital-self-assessment-ipv & No Thailand-specific RCT or \seqsplit{effectiveness} (as opposed to \seqsplit{preliminary} expert-evaluation) data were found for a scored self-check tool; Chuemchit 2026 (PMID 42196762) is \seqsplit{development}\slash{}heuristic-evaluation stage only, not yet tested with the target population or measured for \seqsplit{behavioural}\slash{}health outcomes.\\
\\[3pt]
Thai \slash{} Thai-Muslim samples & humility-compassion-self-regulation & No Muslim-sample or Thailand-based record was found for the \seqsplit{intellectual}-humility or self-compassion \seqsplit{literatures} (angles a\slash{}c); the \seqsplit{religiosity}-IPV evidence located (36331229, 33704630) is Brazilian-mixed-sample and Christian-family literature \seqsplit{respectively}, not Muslim samples, despite the query explicitly requesting Muslim samples where indexed --- none were indexed for this specific construct pairing.\\
\\[3pt]
Thai \slash{} Thai-Muslim samples & ipv-screening-response & No PubMed-indexed evidence located on IPV screening\slash{}response \seqsplit{specifically} among Muslim or Thai\slash{}Southeast Asian couples; the Nigeria trials (Sokoto, Ebonyi) are the closest Muslim-population-adjacent evidence but are clinical, not premarital-education, settings and not Southeast Asian.\\
\\[3pt]
Thai \slash{} Thai-Muslim samples & islamic-medical-ethics-tier-one & No PubMed-indexed record \seqsplit{specifically} evaluates premarital health education content (not screening, not marriage-age demography) delivered through an explicitly named Islamic ethical frame (faqr\slash{}taqwa\slash{}amanah\slash{}'adl\slash{}shura) in a Thai or Southeast Asian Muslim population --- the closest analogues found (Bangladesh mosque-based diabetes RCT, PMID 41114978; Somali refugee mosque-based trauma RCT, PMID 39186273) are faith-i...\\
\\[3pt]
Thai \slash{} Thai-Muslim samples & islamic-medical-ethics-tier-one & No BMC Medical Ethics systematic review specific to Islamic bioethics was found; the one Islamic-adjacent BMC Medical Ethics record located (PMID 41491519, a scoping review on social egg freezing) was excluded as \seqsplit{insufficiently} on-topic (not focused on a Muslim population or the five named ethical terms).\\
\\[3pt]
Thai \slash{} Thai-Muslim samples & \seqsplit{manipulation}-literacy-\seqsplit{inoculation} & No PubMed-indexed Thailand-resident (as opposed to diaspora) study on \seqsplit{recognition} of emotional\slash{}\seqsplit{psychological} abuse or coercive control was found; the two most population-relevant records (26289454, 26442989) are Canada\slash{}US-diaspora South\slash{}Southeast Asian Muslim and Southeast Asian samples, not Thailand-resident \seqsplit{populations} -- a genuine evidence gap for angle (e) as specified.\\
\\[3pt]
Thai \slash{} Thai-Muslim samples & muslim-thailand-southeast-asia & No PubMed-indexed evaluation study of an Islamic\slash{}faith-based premarital education or marriage-\seqsplit{preparation} course (content, uptake, or outcomes) specific to Thailand, Malaysia, or Indonesia was found; premarital-course literature retrieved was \seqsplit{overwhelmingly} about genetic\slash{}\seqsplit{thalassaemia} premarital screening, not relational or \seqsplit{safeguarding} education.\\
\\[3pt]
Thai \slash{} Thai-Muslim samples & muslim-thailand-southeast-asia & No PubMed-indexed study on secret, \seqsplit{unregistered}, or informal (nikah siri \slash{} nikah tidak dicatatkan-type) religious marriage prevalence, drivers, or \seqsplit{consequences} among Muslim \seqsplit{populations} in Thailand\slash{}Malaysia\slash{}Indonesia was located via these queries.\\
\\[3pt]
Thai \slash{} Thai-Muslim samples & muslim-thailand-southeast-asia & No PubMed-indexed study on Thailand's Deep South (Pattani\slash{}Yala\slash{}Narathiwat) Muslim population \seqsplit{specifically} addressing IPV, forced marriage, or premarital programmes was found; the Thailand-specific IPV\slash{}OSCC literature retrieved (\seqsplit{Grisurapong} 2002, Nagae 2004, Phonkacha 2026) is either general-population or historical, not Muslim-population-specific within Thailand.\\
\\[3pt]
Thai \slash{} Thai-Muslim samples & muslim-thailand-southeast-asia & Evidence on Muslim \seqsplit{populations} in Malaysia and Indonesia specific to premarital decision-control, forced-marriage refusal, or STOP-RULE-type joint-versus-separated counseling protocols was not found in PubMed; most Southeast Asian Muslim marriage literature indexed in PubMed concerns child-marriage prevalence\slash{}\seqsplit{determinants} (largely Indonesia) or HIV \seqsplit{vulnerability} within marriage, not programme design or \seqsplit{safeguardin}...\\
\\[3pt]
Thai \slash{} Thai-Muslim samples & muslim-thailand-southeast-asia & Rohingya evidence, while ethnically\slash{}\seqsplit{religiously} proximate and \seqsplit{geographically} linked to the Bay of Bengal\slash{}Southeast Asia border region, is drawn from a Bangladesh refugee-camp setting, not Thailand proper; it is included as the closest available regional Muslim-population IPV\slash{}child-marriage evidence but should not be read as directly \seqsplit{representative} of Thai Muslim or \seqsplit{intercultural} couples.\\
\\[3pt]
Thai \slash{} Thai-Muslim samples & peer-online-support-digital-safe-space & No PubMed-indexed record was found specific to Thailand, or to a Thai Muslim population, on digital peer support, online IPV support, or moderated online \seqsplit{communities}; the Muslim-specific record found (PMID 35536616) is set in Australia, and other Muslim\slash{}health hits from the search (PMID 37928060 anti-Muslim \seqsplit{discrimination}, PMID 38040423 Sikh\slash{}Muslim kidney transplant views, PMID 31120322 US Muslim \seqsplit{contraceptive} bel...\\
\\[3pt]
Thai \slash{} Thai-Muslim samples & premarital-education-outcomes & No Thailand-specific or Southeast Asian premarital\slash{}marriage education outcome study was located; the closest cultural analogues found are Iranian (PMID 40783649), Turkish (PMID 42442188), Saudi (PMID 36620784) and South Korean (PMID 39757299) studies -- none are Southeast Asian or \seqsplit{specifically} Thai Muslim\slash{}\seqsplit{intercultural} \seqsplit{populations}.\\
\\[3pt]
Thai \slash{} Thai-Muslim samples & primary-prevention-frameworks & No papers were found \seqsplit{specifically} evaluating IPV-prevention programme mechanisms among Muslim or \seqsplit{intercultural} couples in Thailand or Southeast Asia; the Nigerian Muslim\slash{}Christian \seqsplit{congregational} trial (PMID 37329160) is the nearest cross-cultural, mixed-faith analogue found.\\
\\[3pt]
Thai \slash{} Thai-Muslim samples & primary-prevention-frameworks & Faith-based \seqsplit{intervention} evidence located skews toward Christian\slash{}Catholic and one Korean-immigrant Christian clergy sample plus one Nigerian Muslim\slash{}Christian trial; \seqsplit{specifically} Islamic religious-authority-led premarital education \seqsplit{evaluations} were not located in this search.\\
\\[3pt]
Premarital-specific trials & assessor-\seqsplit{independence}-checkin-models & No PubMed-indexed literature was found that directly studies 'conflict of interest of a wedding-service provider acting as premarital-education assessor' --- this exact \seqsplit{configuration} does not appear to be a studied topic; evidence is \seqsplit{necessarily} analogical (organ-donor advocacy, clinical dual-\seqsplit{relationship} ethics, facial-plastic-surgery COI) rather than domain-matched.\\
\\[3pt]
Premarital-specific trials & couple-work-safety-coercive-control & Coercive-control \seqsplit{measurement} \seqsplit{instruments} identified (CIPR, C-CAS, LAS, C-CBS-R, CASR-SF) are validated \seqsplit{predominantly} in Western (US, Australian, Canadian) or Chinese-community samples; only one (Tiwari 2015, Chinese sample) offers a non-Western cross-cultural precedent, and none has been validated for premarital screening \seqsplit{specifically} (all are post-hoc\slash{}\seqsplit{retrospective} abuse measures, not \seqsplit{prospective} risk-of-future-c...\\
\\[3pt]
Premarital-specific trials & couple-work-safety-coercive-control & No cost-\seqsplit{effectiveness}, \seqsplit{implementation}-fidelity, or long-term (multi-year) outcome data were found for any couple-based IPV-adjacent \seqsplit{intervention} in this angle's search; the strongest RCT (Taft 2024) is a targeted secondary-prevention trial in a screened military population, not a universal premarital population.\\
\\[3pt]
Premarital-specific trials & decision-autonomy-measures & No study evaluating premarital or marriage-\seqsplit{preparation} education \seqsplit{specifically} through a decision-control or \seqsplit{reproductive}-autonomy outcome measure (rather than knowledge\slash{}attitude\slash{}\seqsplit{satisfaction}) was identified; the ARCHES trial (Miller 2016) is the closest analogue but is a clinic-based, not marriage-\seqsplit{preparation}, \seqsplit{intervention}.\\
\\[3pt]
Premarital-specific trials & decision-autonomy-measures & No study combining religious\slash{}faith-based marriage \seqsplit{preparation} \seqsplit{programming} with a validated decision-autonomy or coercion-\seqsplit{recognition} outcome measure was found in PubMed; such programmes, if evaluated at all, appear to be published outside PubMed-indexed venues or use non-\seqsplit{standardized} outcome measures.\\
\\[3pt]
Premarital-specific trials & digital-self-assessment-ipv & No PubMed-indexed record found \seqsplit{specifically} combining 'premarital' status (not-yet-married couples) with a scored anonymous self-check tool - all identified trials enrol women already \seqsplit{identifying} current\slash{}recent partner abuse, not a general premarital population screening before any disclosed problem exists.\\
\\[3pt]
Premarital-specific trials & ipv-screening-response & No PubMed-indexed evidence located specific to premarital\slash{}pre-nikah education settings combined with IPV screening\slash{}response protocols -- all direct LIVES\slash{}ARCHES \seqsplit{effectiveness} evidence (PMIDs 41267027, 41151831, 32460786) comes from antenatal\slash{}family-planning clinical encounters, not premarital or religious pre-marriage \seqsplit{counselling} contexts. This is a genuine evidence gap for T-PHE's specific setting, not merely a s...\\
\\[3pt]
Premarital-specific trials & islamic-medical-ethics-tier-one & No PubMed-indexed FIMA (Fiqh Council\slash{}Federation of Islamic Medical \seqsplit{Associations}) or IMANA formal position\slash{}guideline paper on premarital counseling, marriage \seqsplit{preparation}, or \seqsplit{reproductive}-decision consent was found (Q7 'FIMA\slash{}IMANA + guideline' returned 0 hits); the two Journal of IMA articles included (PMID 23610513, 23864762) are topical essays, not formal guidelines.\\
\\[3pt]
Premarital-specific trials & \seqsplit{manipulation}-literacy-\seqsplit{inoculation} & No PubMed-indexed record found that directly measures '\seqsplit{recognition} of coercive control\slash{}\seqsplit{gaslighting}' via a validated \seqsplit{psychometric} instrument \seqsplit{specifically} among \seqsplit{adolescents}\slash{}young adults in a premarital-education context (searches for 'coercive control \seqsplit{recognition} scale' and '\seqsplit{gaslighting} \seqsplit{recognition} education' returned 0 or near-0 direct hits; adjacent constructs like IPV \seqsplit{internalized} stigma (PMID 39536684) and gasli...\\
\\[3pt]
Premarital-specific trials & \seqsplit{manipulation}-literacy-\seqsplit{inoculation} & No PubMed-indexed study was found that directly tests whether \seqsplit{recognising}\slash{}labelling a \seqsplit{relationship} pattern as 'abuse' causally predicts subsequent help-seeking in a premarital or pre-\seqsplit{relationship} (as opposed to already-victimized) population -- the help-seeking records found (37882155, 38336660, 26442989) all study people who already \seqsplit{experienced} IPV, not people using an anonymous early-signal self-check before har...\\
\\[3pt]
Premarital-specific trials & peer-online-support-digital-safe-space & No PubMed-indexed record evaluates a premarital-education-linked or nikah\slash{}marriage-\seqsplit{preparation} digital peer community anywhere; this proposal's Layer 3 model appears untested in the indexed literature.\\
\\[3pt]
Premarital-specific trials & peer-online-support-digital-safe-space & No LMIC\slash{}Southeast Asia-specific systematic review of moderated digital peer support for \seqsplit{relationship} or premarital health topics was found via this angle; the closest LMIC-adjacent records (PMID 37540713, 38558241, 39116286) concern general digital health promotion, not peer community moderation \seqsplit{specifically}, and were excluded as off-angle.\\
\\[3pt]
Premarital-specific trials & premarital-education-outcomes & No PubMed-indexed study was found that evaluates a premarital or marriage education programme against decision-control, coercion-\seqsplit{recognition}, or informed-choice outcomes as primary endpoints (the specific outcome domain the FAILURE CRITERION targets) -- every \seqsplit{effectiveness} study located, including the Hawkins\slash{}ACF meta-analyses, measures \seqsplit{satisfaction}, \seqsplit{communication} skills, \seqsplit{relationship} stability, or quality of life...\\
\\[3pt]
Premarital-specific trials & premarital-education-outcomes & No PubMed-indexed evaluation of a STOP RULE-type screen-and-defer mechanism for suspected coercion\slash{}violence within premarital (as opposed to couple\slash{}marital) education \seqsplit{specifically} was found; the strongest supporting evidence (Markman 2013 PMID 23421844; Williamson 2015 PMID 25664643) is indirect, showing that joint \seqsplit{relationship} education can worsen outcomes for high-aggression\slash{}high-risk couples, not a direct test ...\\
\\[3pt]
Premarital-specific trials & premarital-education-outcomes & No PubMed-indexed premarital-education outcome study directly measures 'referral timeliness, \seqsplit{accessibility}, and follow-through' (P5) as a composite outcome; Williamson 2019 (PMID 30489129) and Williamson 2014 (PMID 24294929) are the closest available evidence, addressing help-seeking barriers and gateway effects \seqsplit{respectively}, but neither evaluates an actual referral pathway or its quality.\\
\\[3pt]
Premarital-specific trials & premarital-education-outcomes & No meta-analysis or systematic review \seqsplit{specifically} on faith-based\slash{}religious-leader-delivered premarital education and safety\slash{}coercion outcomes was found via the systematic-review-filtered queries (query 2 returned 0 hits on the exact phrasing tried); this may reflect true absence of such reviews in PubMed rather than a search-term failure, but should not be treated as confirmed absence without a broader grey-liter...\\
\\[3pt]
Premarital-specific trials & primary-prevention-frameworks & No PubMed-indexed paper was found that evaluates a premarital or pre-nuptial education programme \seqsplit{specifically} structured around a governance\slash{}decision-control framework comparable to T-PHE; the closest analogues (SASA!, \seqsplit{Indashyikirwa}, MFP) are ongoing-couple or community-level \seqsplit{interventions}, not marriage-transition-timed \seqsplit{interventions}.\\
\\[3pt]
Premarital-specific trials & primary-prevention-frameworks & Cost-\seqsplit{effectiveness} and \seqsplit{implementation}-science evidence located pertains to community \seqsplit{mobilisation} and CHW-delivered parenting\slash{}couples programmes, not premarital education delivered through religious \seqsplit{institutions}, limiting direct \seqsplit{transferability} to T-PHE's delivery channel.\\
\\[3pt]
Decision-control \seqsplit{instruments} & couple-work-safety-coercive-control & No PubMed-indexed study directly tests a 'stop rule' construct (halt joint work pending \seqsplit{confidential} assessment, never resume if coercive control is confirmed) as an explicit \seqsplit{intervention} design feature; the closest analogues (McCollum 2008, Schneider 2014, Karakurt 2016) describe screening\slash{}safety-\seqsplit{modification} practice \seqsplit{recommendations} rather than a validated stop\slash{}resume protocol.\\
\\[3pt]
Decision-control \seqsplit{instruments} & couple-work-safety-coercive-control & Evidence on screening\slash{}assessing male \seqsplit{perpetration} (rather than \seqsplit{victimization}) is \seqsplit{comparatively} sparse (Portnoy 2020), which matters for T-PHE's \seqsplit{bidirectional}\slash{}decision-control framing (P1) where either partner could be at risk.\\
\\[3pt]
Decision-control \seqsplit{instruments} & decision-autonomy-measures & No PubMed-indexed instrument was found that measures 'decision control' \seqsplit{specifically} at the point of a marriage-transition decision (proceed\slash{}delay\slash{}refuse marriage) for adult couples; the closest analogues are \seqsplit{reproductive}-decision \seqsplit{instruments} (RAS, RCS, FP-ADM) or an adolescent \seqsplit{empowerment} subscale ('choice of partners, marriage, and children' in Upadhyay 2021), none validated for the marriage-entry decision itself.\\
\\[3pt]
Decision-control \seqsplit{instruments} & decision-autonomy-measures & No instrument \seqsplit{specifically} measuring 'coercive control \seqsplit{recognition}' as a discrete competency (as opposed to experience\slash{}\seqsplit{victimization} scales like the RCS) was found; the T-PHE STOP RULE presumes \seqsplit{practitioners} and\slash{}or \seqsplit{participants} can recognize coercive control, but a dedicated \seqsplit{recognition}-literacy instrument (as distinct from an experience-report scale) was not located in this search.\\
\\[3pt]
Decision-control \seqsplit{instruments} & digital-self-assessment-ipv & No record was found addressing a combined single self-check instrument spanning fear, control, pressure, \seqsplit{reproductive} health AND mental health \seqsplit{simultaneously} (as T-PHE's Layer 0 proposes) - the literature covers these domains via separate, purpose-specific tools (IPV decision aids vs Danger Assessment vs \seqsplit{reproductive}-coercion measures), not one integrated scored instrument.\\
\\[3pt]
Decision-control \seqsplit{instruments} & ipv-screening-response & Evidence on harms of screening (adverse effects, loss of privacy, distress) is present (PMID 22565034, 26200817) but almost entirely measured only \seqsplit{immediately} post-screening, with very limited longer-term harm \seqsplit{surveillance} across the literature -- relevant to T-PHE's dignity\slash{}consent safeguard but not fully answered by current evidence.\\
\\[3pt]
Layer 0 \slash{} Layer 2 evidence & digital-self-assessment-ipv & Safety-harm data (partner discovery, \seqsplit{retaliation}) are addressed \seqsplit{qualitatively}\slash{}\seqsplit{methodologically} (Tarzia 2017) but no record reports a quantified adverse-event rate from an anonymous online self-check \seqsplit{specifically}, which the T-PHE failure criterion would need.\\
\\[3pt]
Layer 0 \slash{} Layer 2 evidence & peer-online-support-digital-safe-space & Evidence on moderator training, competence, and error rates in routing acute-risk \seqsplit{disclosures} (self-harm, IPV, abuse) to services is sparse (PMID 34193497 found only 5 usable studies after screening 397); no record quantifies false-negative or delayed-routing rates for non-clinical moderators.\\
\\[3pt]
Layer 0 \slash{} Layer 2 evidence & peer-online-support-digital-safe-space & Evidence \seqsplit{quantifying} how often anonymity itself is defeated by partner technology-\seqsplit{facilitated} \seqsplit{surveillance} of a specific digital peer-support platform (as opposed to general social media\slash{}apps) was not found; existing TFA evidence (PMID 35537445, 32998673, 40271872) describes general digital abuse mechanisms, not platform-specific breach of anonymous peer-support spaces.\\
\\[3pt]
\end{longtable}
